% ---------------------------------------------------------------------
% Parte 6 · Panel de Mario · expone Mario · 1:30 (lo edita solo Mario)
% La captura es del panel funcionando en tu proyecto, tomada por ti y con
% la fecha visible: guárdala como figuras/panel-mario.png.
% Las tres tarjetas tienen la misma estructura en los tres paneles para
% que la comparación de la parte 8 salga sola.
% ---------------------------------------------------------------------
\begin{frame}{Panel de Mario}{Pre-sustentaciones: la coordinación sigue las solicitudes sin recargar}
  \begin{columns}[T]
    \begin{column}{0.56\textwidth}
      \figura{figuras/panel-mario.png}{\linewidth}{3.3cm}{Captura del panel funcionando, con la fecha
        visible (evidencia de primera mano).}
      \begin{tarjeta}{bcmAcento}{Si se cae la conexión}
        \footnotesize Aviso en ámbar y reintento a 1, 2, 4, 8 s, luego cada 10\par
        Al volver, recarga todas las solicitudes
      \end{tarjeta}
    \end{column}
    \begin{column}{0.4\textwidth}
      \begin{tarjeta}{bcmTinta}{Qué hace}
        \footnotesize Muestra cada solicitud con su estado\par
        y la lista de eventos recibidos
      \end{tarjeta}
      \begin{tarjeta}{bcmWS}{Dónde entra el WebSocket}
        \footnotesize El director cambia un estado\par
        y el panel se entera sin preguntar
      \end{tarjeta}
    \end{column}
  \end{columns}
\end{frame}
